% Zone „Das kennst du schon“ – aus bank/grenzwerte-und-verhalten-im-unendlichen/zone.jsonl (zusammenbau v0.1)
\einheitenkopf[zone][Kennst du schon]{Das kennst du schon}
\setcounter{aufgabe}{0}

%% TODO Titel: Ich-kann-Satz fehlt in der Bank (Platzhalter: Kettenname) – Nr. 1
%% TODO Anweisung: ein Satz über der ersten Teilaufgabe fehlt in der Bank – Nr. 1
\begin{aufgabe}{Potenzen mit geradem und ungeradem Exponenten bei negativer Basis}
\begin{teile}
\teil Berechne: $(-2)^4$ \leerfeld
\teil Berechne: $(-0{,}5)^3$ \leerfeld
\end{teile}
\end{aufgabe}

%% TODO Titel: Ich-kann-Satz fehlt in der Bank (Platzhalter: Kettenname) – Nr. 2
%% TODO Anweisung: ein Satz über der ersten Teilaufgabe fehlt in der Bank – Nr. 2
\begin{aufgabe}{Öffnung einer Parabel und Verlauf der Normalparabel dritten Grades nach oben und unten}
\begin{teile}
\teil Wie ist die Parabel $y = 3x^2 - 1$ geöffnet? \kreuz{nach oben} \\ \kreuz{nach unten}
\teil Wie verläuft $y = -x^3$? \kreuz{von links unten nach rechts oben} \\ \kreuz{von links oben nach rechts unten}
\end{teile}
\end{aufgabe}

%% TODO Titel: Ich-kann-Satz fehlt in der Bank (Platzhalter: Kettenname) – Nr. 3
%% TODO Anweisung: ein Satz über der ersten Teilaufgabe fehlt in der Bank – Nr. 3
\begin{aufgabe}{Exponentialfunktion y = a · qˣ}
\begin{teile}
\teil $y = 2^x$: Werte für $x = 1$, $x = 2$, $x = 3$? \leerfeld , \leerfeld , \leerfeld
\teil $y = 3 \cdot 0{,}8^x$: Wert für $x = 2$? y = \leerfeld
\end{teile}
\end{aufgabe}

%% TODO Titel: Ich-kann-Satz fehlt in der Bank (Platzhalter: Kettenname) – Nr. 4
%% TODO Anweisung: ein Satz über der ersten Teilaufgabe fehlt in der Bank – Nr. 4
\begin{aufgabe}{Achsen- und Punktsymmetrie am Graphen erkennen und f(−x) bilden (Vorzeichen bei geraden und ungeraden Potenzen)}
\begin{teile}
\teil $f(x) = x^2 + 3$: Berechne $f(2)$ und $f(-2)$. f(2) = \leerfeld , f(−2) = \leerfeld
\teil $f(x) = 2x^4 - x^2$: Bilde $f(-x)$. f(−x) = \leerfeld
\end{teile}
\end{aufgabe}

%% TODO Titel: Ich-kann-Satz fehlt in der Bank (Platzhalter: Kettenname) – Nr. 5
%% TODO Anweisung: ein Satz über der ersten Teilaufgabe fehlt in der Bank – Nr. 5
\begin{aufgabe}{Satz vom Nullprodukt und die Regel, dass $e^x$ nie null und immer positiv ist}
\begin{gleichungsraster}[2]
\gl{\text{Löse: $(x - 4)(x + 1) = 0$}} &
\gl{\text{Löse: $(2x - 3)(x + 5) = 0$}} \\
\end{gleichungsraster}
\end{aufgabe}

%% TODO Titel: Ich-kann-Satz fehlt in der Bank (Platzhalter: Kettenname) – Nr. 6
%% TODO Anweisung: ein Satz über der ersten Teilaufgabe fehlt in der Bank – Nr. 6
\begin{aufgabe}{Verschiebung und Streckung eines Graphen an der Gleichung lesen (f(x) + c verschiebt nach oben, a · f(x) streckt, negatives a spiegelt) und Monotonie als „steigt“ oder „fällt“}
\begin{teile}
\teil $y = x^2 + 5$: Um wie viel und in welche Richtung ist die Normalparabel verschoben? um \leerfeld nach \leerfeld
\teil Steigt oder fällt $y = 0{,}5 \cdot 2^x$?
\end{teile}
\end{aufgabe}

%% TODO Titel: Ich-kann-Satz fehlt in der Bank (Platzhalter: Kettenname) – Nr. 7
%% TODO Anweisung: ein Satz über der ersten Teilaufgabe fehlt in der Bank – Nr. 7
\begin{aufgabe}{Vorzeichen eines Produkts und eines Bruchs aus den Vorzeichen der Faktoren bestimmen (plus mal minus, Zähler positiv durch Nenner positiv)}
\begin{teile}
\teil Vorzeichen von $(-4) \cdot 7$? \kreuz{positiv} \\ \kreuz{negativ}
\teil Vorzeichen von $(x - 5) \cdot e^{x}$ für $x = 2$? \kreuz{positiv} \\ \kreuz{negativ}
\end{teile}
\end{aufgabe}

%% TODO Titel: Ich-kann-Satz fehlt in der Bank (Platzhalter: Kettenname) – Nr. 8
%% TODO Anweisung: ein Satz über der ersten Teilaufgabe fehlt in der Bank – Nr. 8
\begin{aufgabe}{Vorzeichen eines Produkts und eines Bruchs aus den Vorzeichen der Faktoren bestimmen (plus mal minus, Zähler positiv durch Nenner positiv) – Fehler finden}
\begin{teile}
\teil Ein Schüler schreibt: „$(6 - x) \cdot e^{x}$ ist für $x = 10$ positiv, weil $e^{10}$ positiv ist.“ Finde den Fehler.
\end{teile}
\end{aufgabe}

%% TODO Titel: Ich-kann-Satz fehlt in der Bank (Platzhalter: Kettenname) – Nr. 9
%% TODO Anweisung: ein Satz über der ersten Teilaufgabe fehlt in der Bank – Nr. 9
\begin{aufgabe}{Vorzeichen eines Produkts und eines Bruchs aus den Vorzeichen der Faktoren bestimmen (plus mal minus, Zähler positiv durch Nenner positiv) – selbst rechnen}
\begin{teile}
\teil Vorzeichen von $(3 - x) \cdot e^{2x}$ für $x = 8$? Begründe mit den Faktoren.
\end{teile}
\end{aufgabe}
